\documentclass[../../../include/open-logic-section]{subfiles}

\begin{document}
	
\olfileid[tr]{sth}{replacement}{finiteaxiomatizability}
\olsection{Ek: Sonlu Aksiyomlaştırılabilirlik}
Bu bölümü Yerine Koyma'dan bazı sonuçlar çıkararak kapatıyoruz. İlk sonuç
\citet{Montague1961}'e aittir; bunun $\ZF$ \emph{içinde} bir ispat değil,
$\ZF$ \emph{hakkında} bir ispat olduğuna dikkat edin:
\begin{thm}\ollabel{zfnotfinitely}
	$\ZF$ sonlu olarak aksiyomlaştırılamaz. Daha genel olarak: $\Th{T}$ sonlu
	ve $\Th{T} \Proves \ZF$ ise $\Th{T}$ tutarsızdır.
	
	(Burada kendimizi, mantık dışı tek ilkel simgesi $\in$ olan birinci
	dereceden tümcelerle örtük olarak sınırlandırıyor ve bütün $\phi \in \ZF$
	için $\Th{T} \Proves \phi$ olduğunu belirtmek üzere
	$\Th{T} \Proves \ZF$ yazıyoruz.)
\end{thm}
\begin{proof}
	$\Th{T} \Proves \ZF$ olacak biçimde sonlu bir $\Th{T}$ sabit olsun.
	Dolayısıyla $\Th{T}$ Yansıma'yı, yani
	\olref[sth][replacement][ref]{reflectionschema} sonucunu ispatlar.
	$\Th{T}$ sonlu olduğundan onu tek bir $\theta$ evetlemesi olarak
	yeniden yazabiliriz. Bu formülle yansıma yapınca $\Th{T} \Proves \exists
	\beta(\theta \liff \theta^{V_\beta})$ olur. Önemsiz biçimde
	$\Th{T} \Proves \theta$ olduğundan $\Th{T} \Proves \exists \beta\
	\theta^{V_\beta}$ sonucunu elde ederiz.
	
	Şimdi $\psi(X)$ şunu kısaltsın:
	\[
		\theta^X \land X\text{ geçişlidir} \land (\forall Y \in X)(Y\text{ geçişlidir}\lif \lnot \theta^{Y})
	\]
	Bu kabaca $X$'in geçişli bir $\theta$ modeli ve bu bakımdan $\in$-minimal
	olduğunu söyler. $\Th{T} \Proves \exists \beta\ \theta^{V_\beta}$
	olduğunu hatırlayıp $\ZF$ içindeki ve dolayısıyla $\Th{T}$ içindeki
	ranklara ilişkin temel olguları kullanınca şunu elde ederiz:
	\begin{equation}
		\Th{T} \Proves \exists M \psi(M). \tag{*}\label{Mpsi}
	\end{equation}
	$\psi(X)$'in ilk evetlenenini kullanırsak, $\Th{T} \Proves \sigma$
	olduğu her durumda $\Th{T} \Proves \forall X(\psi(X) \lif \sigma^X)$
	elde ederiz. Dolayısıyla \eqref{Mpsi} gereği:
	\begin{align*}
		\Th{T} &\Proves \forall X(\psi(X) \lif (\exists N \psi(N))^X)\\
	\intertext{Bunu ve yine \eqref{Mpsi} sonucunu kullanırsak:}
		\Th{T} &\Proves \exists M(\psi(M) \land (\exists N \psi(N))^M)
	\intertext{O hâlde özellikle:}
		\Th{T} &\Proves \exists M(\psi(M) \land (\exists N \in M)((N\text{ geçişlidir})^N \land (\theta^N)^M))
	\intertext{Dolayısıyla geçişlilik hakkındaki temel akıl yürütmeyle:}
		\Th{T} &\Proves \exists M(\psi(M) \land (\exists N \in M)(N\text{ geçişlidir} \land \theta^N))
	\end{align*} 
	Böylece $\Th{T}$ tutarsızdır.\footnote{Bu ``temel akıl yürütme'', geçişli
	kümeler için bazı ``mutlaklık olgularını'' ispatlamayı içerir.}
\end{proof}

\citet[223]{Potter2004} tarafından belirtilen benzer bir sonuç şöyledir:

\begin{prop}\ollabel{finiteextensionofZ}
	$\Th{T}$, $\Z$'yi sonlu sayıda yeni aksiyomla genişletsin.
	$\Th{T} \Proves \ZF$ ise $\Th{T}$ tutarsızdır. (Burada
	\olref{zfnotfinitely} için kullandığımız örtük sınırlamaların aynısını
	kullanıyoruz.)
\end{prop}
\begin{proof}
	Ayırma'nın (sonsuz sayıdaki) örnekleri \emph{dışında} $\Th{T}$'nin bütün
	aksiyomlarının evetlemesi için $\theta$ kullanın. $\psi$'yi
	\olref{zfnotfinitely} içindeki gibi $\theta$'dan tanımlayarak
	$\Th{T} \Proves \exists M \psi(M)$ olduğunu gösterebiliriz.
	
	\olref{zfnotfinitely} içinde olduğu gibi, $\Th{T} \Proves \sigma$ olduğu
	her durumda $\Th{T} \Proves \forall X(\psi(X) \lif \sigma^X)$ olduğunu
	veren şemayı kurabiliriz. Ardından ispatımızı tam olarak
	\olref{zfnotfinitely} içindeki gibi tamamlarız.
	
	Ancak şemayı kurmak \olref{zfnotfinitely} içindekinden biraz daha fazla
	çalışma gerektirir. Sonuçta Ayırma örnekleri $\Th{T}$ içindedir, ama
	$\theta$'nın evetlenenleri değildir. Yine de her Ayırma örneği
	$\sigma$ için $\Th{T} \Proves \forall X(X\text{ geçişlidir} \lif
	\sigma^X)$ olduğunu ispatlayarak bu engeli aşabiliriz. Bunu okura
	bırakıyoruz.
\end{proof}
\begin{prob}
	Her Ayırma örneği $\sigma$ için $\Z \Proves \forall X(X\text{ geçişlidir}
	\lif \sigma^X)$ olduğunu gösterin. (Bu şemayı
	\olref[sth][replacement][finiteaxiomatizability]{finiteextensionofZ}
	içinde kullandık.)
\end{prob}
\begin{prob}
	Her $\phi \in \Z$ için $\ZF \Proves \phi^{V_{\omega+\omega}}$ olduğunu
	gösterin.
\end{prob}
\begin{prob}
	\olref[sth][replacement][finiteaxiomatizability]{zfnotfinitely} ve
	\olref[sth][replacement][finiteaxiomatizability]{finiteextensionofZ}
	ispatlarında kullanılan geri kalan şematik sonuçları doğrulayın.
\end{prob}

\olref[sth][replacement][absinf]{sec} içinde belirtildiği gibi bu, Yerine
Koyma'nın \olref[ordinals][ordtype]{thmOrdinalRepresentation} sonucundan
kesin olarak daha güçlü olduğunu gösterir. Ya da biraz daha kesin söylersek:
$\Z$ + ``her iyi sıralama biricik bir sıra sayısına izomorftur'' tutarlıysa
bazı Yerine Koyma örneklerini ispatlayamaz.


	% Varsayım gereği, $\psi^{V_\alpha}$ şu biçimdedir: 
	% \[
	% 	(\forall A \in V_\alpha)(\exists S \in V_\alpha)(\forall x \in V_\alpha)(x \in S \liff (\phi^{V_\alpha}(x) \land x \in A))
	% \]
	% Bunun geçerli olduğunu kurmak için $A \in V_\alpha$ sabit olsun. Ayırma ile şunu elde edin:
	% \[
	% 	S = \Setabs{x \in A}{\phi^{V_\alpha}(x)}
	% \]
	% Şimdi $S \subseteq A \in V_\alpha$ olduğundan $S \in V_\alpha$ ve açıkça $(\forall x \in V_\alpha)(x \in S \liff (\phi^{V_\alpha}(x) \land x \in A)$ olur.

\end{document}
